%% Religion and other complex stuff is stored here
%% It will get an own book

%% TODO review

\section {How to write your own adventures}

Solarpunk adventures are a bit different than other TTRPG adventures - they should transport a specific positive mindset. You are still not limited ! Here are my experiences so far. Please share your insights as well.

Traditional adventures will have a showdown and some minor fights as challenges to create action. In a Solarpunk adventure you will very likely have no fights at all - or use them in a different way. In the "World destroying machine" the boss fight at the end just causes much more trouble. The better way ist to just anaesthetise the hamster. The protagonists also do not start with weapons. If they feel like they need some, they will have to build them.

Better challenges are social challenges, solving dilemmas and building communities.

\subsection{Produce, not consuming}

The Solarpunk world is a world where passive consumers don't have a place. Even the Norms - where the basic needs are covered - tend to accumulate hobbies, arts and social projects instead of becoming passive. If you introduce Solarpunky NPCs in your adventures be prepared to answer: what do they contribute to society and the world ? How ? And Why ?

\subsection{Protagonists}

The player characters are the protagonists. Not necessarily the heroes. While these protagonists are active, the other people in this world are trying to make it a better place as well. So it would be a valid Solarpunk adventure if an NPC group is the movers and shakers and will save the day - but the help from the protagonists is needed to get them unstuck. Maybe create an adventure where they are the sidekicks for the real heroes for a change.

\subsection{Green zones}

Your adventure will very likely need situations and locations where things need fixing. The Solarpunk 2050 setting is in the Solarpunk revolution phase, in some places the situation is closer to the "dirty road to eden" phase than to the "established Solarpunk" phase. The players and their protagonists still deserve a fuzzy feeling after finishing an adventure. So plan for "Green zones". Either positive Solarpunk areas where the adventure starts and the victory after the adventure will be celebrated. Or their success solving the trouble is so incredible that they just created a Green zone. Celebrate there.

Close to the almost perfect green zones are awesome areas where people experimented with Solarpunk style solutions to solve their local problem. The inhabitants solved their specific problems in their own way. This solution might not be perfect, maybe it is not "one size fits all". But it is the solution that worked for the people living there. And it might be the wrong approach for other factions or whatever. For most visitors it will feel strange there.

Those locations can be a hook for a story - and the reason why the protagonists are called.

Introducing and experimenting with unusual concepts how to clean up after a disaster are a bonus here.

\subsection{Real world}

Solarpunk 2050 is set in a realistic world. The tech level is not absurd. And the locations are taken from real towns, cities, landscapes and areas. A map of locations already covered can be found here: \href{https://solarpunk2050.com/map.php}{https://solarpunk2050.com/map.php}.

I would love to add your adventure and location to the map.

For brainstorming I am making a short list of oddities and potential disasters that can happen there.

In the case of Ravensburg I read in the newspaper that the city Wifi is switched off during the night to reduce electro smog.
next step: What if I give them more time and the make it more extreme ? Well: They switch off electricity during the night.
step 3: maximize the impact: Lost would not feel any impact at all... so I made Ravensburg a Norm city.
step 4: Adoption. What is the most positive adoption they can find to align their life with this situation ? Low tech entertainment at night, sleeping early, relaxed evenings, candle light dinners

If you have a town, you can also start by pondering: how would the different fractions inhabit this city in 2050 ? Would they blend ? Or stay in separate neighbourhoods ? How would they shape their district ?
Are there tensions ? Maybe not between fractions but inner fraction tensions with disagreement on a specific project ?

\subsection{Fights}

There are less fights in a Solarpunk setting than in a Dungeons and Dragons setting. Even having no fights at all is quite likely. As you are playing this game to have fun you need other ways to create dynamic situations.

My fixes so far:
\begin{itemize}
    \item{Tuning the pulp factor to 11}
    \item {Create action by races against natural disaster or similar dynamics}
    \item{Create conflicts between the three philosophies. Especially conflicts within the group of protagonists}
    \item {Deadliness is at about 3, the danger is the environment}
    \item {Fights are at 1. They are possible. And will always be the solution with the most side effects}
    \item {Awesomeness is set to 11. As soon as the people Solarpunkified their environment in any way it will feel awesome. Maybe strange and it is maybe not really the thing the players would love, but can accept it is awesome for a certain type of people}
    \item {Crazy people set to 8: Lots of unique and diverse people sticking together and saving their home. With the diversity as their asset}
\end{itemize}

Footnote about diversity: In every role playing game the smartest group build is a diverse group. Healer, Fighter, Archer, Mage, ... That way there is always someone with the skills to solve the different challenges.

\subsection{Topic, moral or theme}

A hack for me to keep the adventures unique was to define a main topic/theme/moral at the beginning. Something easy to remember for me as GM and as a fall back as soon as I have to decide in a hurry.
In adventures I wrote I write it at the beginning. It is not as important as you might think - I do not want to lecture anyone. Most of the time I develop it after half of the adventure is already written. So ir really is "just" a guiding light to keep this story a bit different than the others.

Most of those are fetched from the news cycle. The topics that worked the best so far is what the capitalistic systems (does not really matter if western capitalism or China) did to the global south. Science news about new prototypes work great as well. This can be transplanted into a disaster situation. And the resulting in an enhanced disaster situation with a bad "solution". The characters will have a social/political mess to solve and a real physical problem. Fun for everyone.

I think the players never realized there is this small section at the beginning of my adventures. And I never rubbed the moral in.

\subsection{Mixing genres}

One trick for me to create different adventures is to mix in other genres. Many genres worked very well with the Solarpunk setting. This is basically the trick Joss Whedon ran when combining Vampire hunters and High School drama, or Cowboys and Space.

\subsection{Motivation for the protagonists}

As long as the protagonists do not know many NPCs that could need help yet, there are several generic motivations to get them started. Of course: Some Characters may have their own intrinsic motivations to search the adventure.

\begin{itemize}
\item{Getting Resource points}
\item{A Lost character could be paid in rare books}
\item{A Norm character in a special events with their favourite band}
\item {Pioneers could get access to new tools and technology}
\item {The local community needs help - the protagonists are affected directly}
\item {Gaia priests can also be a source for new adventures. Not only for believers}
\item{Collect the coupon adventures - the protagonists need 5 rare books to enter Alexandria. This is the motivation for the next 5 adventures.}
\end{itemize}


\subsection{Technology}

If the adventure introduces new technology - go for it. But please keep in mind that the tech level for pioneers is limited to "things that are in prototype state in the real world". Norms will focus more on mass production/automation and use the older technology - mostly advanced versions of things you can buy today. And Lost will use old technology, but in smart ways, upcycles and repaired.

\subsection{Publishing}

Writing for your own group is a no brainer. But if you want to publish your adventure (free or paid): Go ahead !

But for the players it would be best if the whole world and the adventures there work together without crazy plot holes. So please contact me and we can weave it into the whole fabric together. I can link to it from the page, add the blip to the map. If you wrote a short adventure I can tie it together to an adventure anthology. On Drive-Thru RPG there is a setting to split the earned money between authors.

So all variants are possible. It is important that the players have a convenient access and can combine adventures from different sources.

\subsection{AI}

Don't use it if you publish through my Drive-Thru channels. I have the "No AI used " flag set.

AI is tricky to get right. Read Doctorow on "Centaurs" and "Reverse Centaurs". And the hard part is right now to not be blinded by the marketing. From a Solarpunk perspective there are some additional pros and cons:

\begin{itemize}
\item{Building own skills is a Solarpunk thing}
\item{Using tools is a Solarpunk thing - that would include AI as long as it stays a tool}
\item{Independence from large corps is a Solarpunk thing}
\item{Disconnecting from capitalism and exploitation is a Solarpunk thing}
\item{Embracing new technology is a Solarpunk thing}
\item{Using new technology for the good of the community in a mindful way is a Solarpunk thing}
\item{Not destroying natural resources is a Solarpunk thing}
\end{itemize}

I was not able to navigate this mine field yet. As a Software Engineer I have been building some machine learning systems in the last 20 years to improve my work-life balance and the balance of my coworkers. But the new LLM based AI together with all the marketing smoke and mirrors are really hard to judge. Sorry, no simple solutions for you here.


\subsection{Images}

If your skill is painting: Please join the team !

%% TODO review end
